% Options for packages loaded elsewhere
\PassOptionsToPackage{unicode}{hyperref}
\PassOptionsToPackage{hyphens}{url}
%
\documentclass[
]{article}
\usepackage{amsmath,amssymb}
\usepackage{iftex}
\ifPDFTeX
  \usepackage[T1]{fontenc}
  \usepackage[utf8]{inputenc}
  \usepackage{textcomp} % provide euro and other symbols
\else % if luatex or xetex
  \usepackage{unicode-math} % this also loads fontspec
  \defaultfontfeatures{Scale=MatchLowercase}
  \defaultfontfeatures[\rmfamily]{Ligatures=TeX,Scale=1}
\fi
\usepackage{lmodern}
\ifPDFTeX\else
  % xetex/luatex font selection
\fi
% Use upquote if available, for straight quotes in verbatim environments
\IfFileExists{upquote.sty}{\usepackage{upquote}}{}
\IfFileExists{microtype.sty}{% use microtype if available
  \usepackage[]{microtype}
  \UseMicrotypeSet[protrusion]{basicmath} % disable protrusion for tt fonts
}{}
\makeatletter
\@ifundefined{KOMAClassName}{% if non-KOMA class
  \IfFileExists{parskip.sty}{%
    \usepackage{parskip}
  }{% else
    \setlength{\parindent}{0pt}
    \setlength{\parskip}{6pt plus 2pt minus 1pt}}
}{% if KOMA class
  \KOMAoptions{parskip=half}}
\makeatother
\usepackage{xcolor}
\setlength{\emergencystretch}{3em} % prevent overfull lines
\providecommand{\tightlist}{%
  \setlength{\itemsep}{0pt}\setlength{\parskip}{0pt}}
\setcounter{secnumdepth}{-\maxdimen} % remove section numbering
\ifLuaTeX
  \usepackage{selnolig}  % disable illegal ligatures
\fi
\IfFileExists{bookmark.sty}{\usepackage{bookmark}}{\usepackage{hyperref}}
\IfFileExists{xurl.sty}{\usepackage{xurl}}{} % add URL line breaks if available
\urlstyle{same}
\hypersetup{
  hidelinks,
  pdfcreator={LaTeX via pandoc}}

\author{}
\date{}

\begin{document}

\hypertarget{lezione-4-la-risoluzione}{%
\section{Lezione 4: La risoluzione}\label{lezione-4-la-risoluzione}}

La \textbf{risoluzione} non è altro che la \textbf{densità di pixel}, cioè di puntini, \textbf{che troviamo in una certa area}. Di solito si usa il \textbf{pollice} come \textbf{unità di misura.}

\textbf{Più punti ci sono in quello spazio, più l'immagine sarà ricca di dettagli.} E come la misuriamo, questa risoluzione? Beh, fondamentalmente in due modi. Il primo è la densità dei pixel, i famosi \textbf{DPI}, \textbf{\emph{Dots Per Inch}}, che ci dicono appunto \textbf{quanti puntini ci sono in un pollice}.

L'altro modo è contare tutti i pixel presenti nell'intera immagine, è quello che facciamo quando, per esempio, parliamo di una fotocamera da 16 megapixel.

Il \textbf{megapixel} è un'unità di misura della \textbf{risoluzione} e della qualità di un'apparecchiatura di ripresa, come una fotocamera digitale.

\begin{itemize}
\tightlist
\item
  \textbf{Rapporto numerico:} Per definizione, \textbf{1 Megapixel corrisponde a un milione di pixel}. Ad esempio, un'immagine da 8 Megapixel contiene 8 milioni di pixel.
\item
  \textbf{Qualità del sensore:} Nelle fotocamere, questo parametro indica il numero di celle (sensori elementari) presenti sul circuito di ripresa (CCD).
\item
  \textbf{Impatto visivo:} Più è grande la matrice di sensori (e quindi maggiore è il numero di Megapixel), superiore sarà la qualità e il grado di dettaglio della fotografia, poiché l'immagine sarà composta da un numero più elevato di punti.
\end{itemize}

Questa è la parte affascinante. È quasi un'illusione ottica.

In realtà lo schermo mostra solo una griglia di punti luminosi. Ma se questi punti sono abbastanza piccoli e vicini, il nostro cervello fa il resto, li unisce e ci fa \textbf{percepire linee, curve, forme complesse}. È lo stesso principio di un mosaico, da vicino vedi le singole tessere, ma da lontano vedi un'immagine completa.

E qui le cose si fanno più interessanti. Bisogna capire che la parola risoluzione assume i significati diversi a seconda del contesto. \textbf{Una cosa è catturare un'immagine, un'altra è mostrarla.}

\hypertarget{risoluzione-di-resa-e-di-ripresa}{%
\subsubsection{Risoluzione di resa e di ripresa}\label{risoluzione-di-resa-e-di-ripresa}}

Non sono la stessa cosa. Distinguiamo bene. Da un lato c'è la \textbf{risoluzione di acquisizione o di ripresa.}
È il \textbf{dettaglio che un dispositivo}, come una fotocamera o uno scanner, \textbf{riesce a catturare} dal mondo reale.

Dall'altro abbiamo la \textbf{risoluzione di riproduzione o di resa}, che è il \textbf{dettaglio} che un monitor o una stampante \textbf{riesce a mostrare}. Sono due mondi con regole e necessità completamente diverse.

Prendiamo un monitor standard, per esempio. La sua risoluzione tipica è di circa 72 dpi. È uno standard o storico, che risale ai tempi dei primi computer Apple.

Certo, oggi molti schermi sono più densi, ma questo resta un po' il punto di riferimento per il web. Può sembrare un valore basso, ma per uno schermo, che mette luce direttamente verso i nostri occhi, spesso è più che sufficiente. Ed ecco il contrasto che spiega tutto.

Una stampante può arrivare a superare i 3000 dpi. Ma come è possibile una differenza così enorme? Semplice. Un pixel su uno schermo è un quadratino di luce che cambia colore.

Fatto. Sulla carta, invece, la stampante deve usare centinaia di puntini microscopici di inchiostro di colori diversi, ciano, magenta, giallo, nero, per creare l'illusione di un singolo punto colorato. Ecco perché per la stampa serve una densità di punti infinitamente maggiore per ottenere un'immagine nitida.

E questo finalmente svela perché un'immagine a 72 dpi perfetta a schermo diventa un disastro una volta stampata.

\begin{figure}
\centering
{image.png}
\caption{image.png}
\end{figure}

Adesso è il momento di vedere, con i nostri occhi, cosa significa, in pratica, cambiare la risoluzione di un'immagine. Partiamo da qui. Un'immagine di un rosone a 200x200 pixel.

Osserviamo bene. I contorni sono netti. Si distinguono i singoli pannelli di vetro colorato.

\begin{figure}
\centering
{image.png}
\caption{image.png}
\end{figure}

È un'immagine definita, piena di informazioni visive. E ora dimezziamo la risoluzione. 100x100 pixel.

La differenza è subito evidente, no? I dettagli più piccoli si sono fusi insieme. I bordi sono meno netti. L'immagine ha perso la sua croccantezza.
Stiamo letteralmente buttando via informazioni.

\begin{figure}
\centering
{image.png}
\caption{image.png}
\end{figure}

E infine il colpo di grazia. Crolliamo a 20x20 pixel.

Del rosone non è rimasto quasi nulla. Ora è solo una griglia di quadratoni colorati. Questa è la prova più chiara di tutte.

\begin{figure}
\centering
{image.png}
\caption{image.png}
\end{figure}

I pixel sono l'informazione. \textbf{Meno pixel abbiamo, meno informazione c'è.} Questo confronto è ancora più impietoso.

\begin{figure}
\centering
{image.png}
\caption{image.png}
\end{figure}

A sinistra un'immagine a 1250 dpi. Ogni singolo dettaglio dell'orologio è perfettamente leggibile. A destra la stessa identica immagine, ma a 72 dpi.
Guardiamo come i numeri, le lancette, le incisioni, è diventato tutta una macchia indistinta. La differenza è abissale.
Bene, adesso che abbiamo capito i principi, vediamo come si applicano alle cose che usiamo ogni giorno, dalla carta stampata ai nostri schermi.

Insomma, ogni tipo di stampa ha le sue esigenze. Un quotidiano, con la sua carta porosa che assorbe l'inchiostro, si accontenta di circa 75 dpi. Per una rivista patinata, più lucida e definita, si sale a 133.

Una brochure di qualità ne richiede almeno 175. Fino ad arrivare ai libri fotografici, dove per una resa impeccabile si possono raggiungere risoluzioni altissime, anche 2400 dpi. E adesso, passiamo agli schermi.

\begin{figure}
\centering
{image.png}
\caption{image.png}
\end{figure}

Tutti abbiamo sentito parlare di VGA, HD, 4K, ma cosa significano davvero queste sigle in termini di numero di pixel? Questo grafico può spaventare un po', ma in realtà racconta una storia molto semplice, l'evoluzione della definizione sui nostri schermi. Si parte dai vecchi standard come il VGA, con i suoi 640x480 pixel, poi, passo dopo passo, siamo arrivati all'HD 720, al Full HD o 1080p, che è lo standard 1920x1080 pixel e così via. Ogni gradino di questa scala ha significato più pixel e quindi immagini più nitide.

E con questo arriviamo alla regola d'oro, il concetto più importante da portarsi a casa. Per avere la qualità migliore, un'immagine dovrebbe essere vista alla stessa risoluzione con cui è stata creata. Ingrandire un'immagine a bassa risoluzione su uno schermo ad alta risoluzione non la migliora, anzi, ne mette solo in evidenza tutti i difetti.

\end{document}
